% Resume: Arumalla Mohit Krishna
% Compile with pdfLaTeX.

\documentclass[letterpaper,10pt]{article}

\usepackage[margin=0.5in]{geometry}
\usepackage[T1]{fontenc}
\usepackage{lmodern}
\usepackage{microtype}
\usepackage{titlesec}
\usepackage{enumitem}
\usepackage[hidelinks]{hyperref}
\usepackage{tabularx}
\input{glyphtounicode}

\pdfgentounicode=1
\pagestyle{empty}
\urlstyle{same}
\raggedright
\setlength{\parindent}{0pt}
\setlength{\tabcolsep}{0pt}

\titleformat{\section}
  {\large\scshape}{}{0em}{}[\titlerule]

\titlespacing*{\section}{0pt}{6pt}{3pt}

\newcommand{\entry}[2]{%
  \noindent
  \begin{tabularx}{\linewidth}
    {@{}>{\raggedright\arraybackslash}X@{\hspace{6pt}}r@{}}
    \textbf{#1} & #2
  \end{tabularx}\par
}

\newcommand{\bullets}{
  \begin{itemize}[
    leftmargin=1.2em,
    topsep=2pt,
    itemsep=1pt,
    parsep=0pt,
    partopsep=0pt
  ]
}

\newcommand{\finishbullets}{
  \end{itemize}
}

\begin{document}

%----------HEADING----------
\begin{center}
    {\LARGE\bfseries Arumalla Mohit Krishna}\\[3pt]
    +1 (631) 352-8643 $|$
    \href{mailto:mohitkrishna.arumalla@stonybrook.edu}
    {mohitkrishna.arumalla@stonybrook.edu} $|$
    \href{https://linkedin.com/in/mohit-krishna-arumalla}
    {LinkedIn}
\end{center}
\vspace{-8pt}

%-----------EDUCATION-----------
\section{Education}

\entry
  {\textbf{Stony Brook University}}
  {Aug 2025 -- Dec 2026}

M.S. in Computer Science \hfill New York, USA\\
\textit{
Coursework: Natural Language Processing,
Introduction to Computer Vision, Vision Language Modeling
}
\vspace{3pt}

\entry
  {\textbf{Birla Institute of Technology and Science, Pilani}}
  {Aug 2019 -- 2023}

B.E. in Electronics \& Communication; GPA: 8.02/10 \hfill India
%-----------PROFESSIONAL SUMMARY-----------
\section{Professional Summary}

Software engineer with 2+ years of experience in \textbf{complex distributed systems, AI Harness Engineering, and secure enterprise system integration}, pursuing an M.S. in Computer Science at Stony Brook University. Expertise in \textbf{AI/ML service integration, microservices architecture, cloud infrastructure automation, and DevSecOps practices}. Hands-on experience with \textbf{AI hallucination detection, token optimization, and scalable agentic frameworks} using Python, Go, C++, and infrastructure tools like Docker and Kubernetes.

%-----------TECHNICAL SKILLS-----------
\section{Technical Skills}

\textbf{Languages:}
Python, Go, C++, Java, SQL\\

\textbf{AI \& ML Frameworks:}
PyTorch, TensorFlow/Keras, Transformers, Self-Supervised Learning,
Token Optimization, Anti-Hallucination Detection, BGE, Vision Transformers\\

\textbf{Systems \& Data Engineering:}
Microservices, gRPC, Kafka, Debezium/CDC, PostgreSQL, Redis, Elasticsearch,
Flink SQL, Airflow, BigQuery\\

\textbf{Infrastructure \& DevSecOps:}
Docker, Kubernetes, CI/CD Pipelines, Istio/Envoy, Terraform (basic),
AWS, Linux, New Relic Monitoring

%-----------EXPERIENCE-----------
\section{Experience}

\entry
  {\textbf{Software Engineer II} $|$ \textbf{Infra.Market}, Bangalore}
  {Apr 2024 -- Aug 2025}

\bullets

  \item Built \textbf{semantic product search integrating AI embeddings and vector search}, enabling natural language interaction and retrieval over large product catalogs.

  \item Engineered a CDC pipeline using \textbf{Flink, Kafka, Debezium} ensuring \textbf{exactly-once processing} and reducing latency from minutes to sub-second for real-time data sync across microservices.

  \item \textbf{Deployed and optimized Istio/Envoy service mesh and integrated distributed tracing}, enhancing \textbf{observability and performance of microservices} in production environments.

\finishbullets

\entry
  {\textbf{Software Engineer I} $|$ \textbf{Infra.Market}, Chennai}
  {Jul 2023 -- Mar 2024}

\bullets

  \item Owned Search and Location microservices handling 500K+ daily gRPC requests, ensuring high-availability and \textbf{service resilience} under load.

  \item Remediated PostgreSQL connection leaks via ORM improvements, \textbf{reducing resource exhaustion and improving backend stability by 85\%} during peak traffic.

  \item Developed Airflow pipelines to migrate data from BigQuery to Elasticsearch, including \textbf{dynamic query optimization and automated health monitoring} for search indices.

\finishbullets

%-----------PROJECTS-----------
\section{Projects}

\entry
  {GPU-Optimized Personalized Product Search}
  {Aug 2026 -- Oct 2026}

\textit{Python, PyTorch, BGE, NVIDIA Triton, Go, gRPC}

\bullets

  \item Developed two-tower retrieval model with \texttt{BAAI/bge-small-en-v1.5} embeddings, integrating \textbf{user interaction features} to personalize search relevance.

  \item Built an inference pipeline featuring approximate nearest neighbor retrieval and GPU-accelerated serving using \textbf{NVIDIA Triton}, setup for scalable production workloads.

  \item Benchmarked system under varying loads with detailed latency, throughput, GPU utilization, and cost metrics, validating robustness and \textbf{scalability for enterprise deployment}.

\finishbullets

\entry
  {Attention-Based Hallucination Detection in NMT}
  {Apr 2026 -- May 2026}

\textit{Python, PyTorch, mBART, Transformers, scikit-learn}

\bullets

  \item Created a hallucination detection pipeline for English-German NMT using attention features extracted from \textbf{state-of-the-art transformer models} trained on quality estimation datasets.

  \item Engineered features such as entropy and attention mass ratios across multiple decoder layers and heads to capture hallucination signals for sequence-level and token-level classification.

  \item Achieved strong detection metrics (0.712 F1, 0.695 ROC-AUC), proposing methods for hallucination-aware reranking of translations to improve output reliability in real applications.

\finishbullets

\entry
  {Prompt/LoRA for Self-Supervised Learning}
  {Ongoing}

\textit{Python, PyTorch, Vision Transformers}

\bullets

  \item Researching domain-specific visual prompting and parameter-efficient adaptation (LoRA) in a self-supervised vision transformer framework to enhance \textbf{transfer learning capabilities}.

  \item Evaluating model performance and parameter trade-offs under domain shift for potential applications in adaptive AI model deployment.

\finishbullets

\end{document}